\documentclass[12pt]{article}

%% ── Packages ─────────────────────────────────────────────────────────────────
\usepackage[margin=1.25in]{geometry}
\usepackage{setspace}
\usepackage{amsmath,amssymb}
\usepackage{booktabs}
\usepackage{longtable}
\usepackage{array}
\usepackage{multirow}
\usepackage{tabularx}
\usepackage{graphicx}
\usepackage{float}
\usepackage{caption}
\usepackage{subcaption}
\usepackage{natbib}
\usepackage[colorlinks=true, linkcolor=blue!60!black, citecolor=blue!60!black,
            urlcolor=blue!60!black]{hyperref}
\usepackage[table]{xcolor}
\usepackage[T1]{fontenc}
\usepackage{pdflscape}
\usepackage{rotating}
\usepackage{threeparttable}
\usepackage{dcolumn}
\usepackage{siunitx}
\usepackage{fancyhdr}
\usepackage{titlesec}
\usepackage{abstract}
\usepackage{appendix}

%% ── Layout ───────────────────────────────────────────────────────────────────
\doublespacing
\setlength{\parindent}{0.5in}
\setlength{\parskip}{0pt}

\titleformat{\section}{\large\bfseries}{\thesection.}{0.5em}{}
\titleformat{\subsection}{\normalsize\bfseries}{\thesubsection.}{0.5em}{}

\pagestyle{fancy}\fancyhf{}
\rhead{\small\textit{Communication as Governance}}
\lhead{\small\textit{Working Paper}}
\rfoot{\thepage}
\renewcommand{\headrulewidth}{0.4pt}

\newcommand{\floatfoot}[1]{\par\vspace{3pt}{\footnotesize\textit{#1}}}

%% ── Column types ─────────────────────────────────────────────────────────────
\newcolumntype{L}[1]{>{\raggedright\arraybackslash}p{#1}}
\newcolumntype{C}[1]{>{\centering\arraybackslash}p{#1}}
\newcolumntype{R}[1]{>{\raggedleft\arraybackslash}p{#1}}

%% ── Custom commands ──────────────────────────────────────────────────────────
\newcommand{\PCI}{\textsc{pci}}
\newcommand{\IP}{\textsc{ip}}
\newcommand{\TTO}{\textsc{tto}}
\newcommand{\BERT}{\textsc{bert}}

%% ─────────────────────────────────────────────────────────────────────────────
\begin{document}
\thispagestyle{empty}

\begin{center}
{\LARGE\bfseries Communication as Governance:\\[6pt]
Evidence from University Intellectual Property Policies}

\vspace{20pt}
{\normalsize Hoover Institution Working Paper}\\[4pt]
{\normalsize \today}
\vspace{24pt}
\end{center}

%% ── Abstract ─────────────────────────────────────────────────────────────────
\begin{abstract}
\noindent
Formal policies govern behavior not only through the substance of rules but through how those rules are communicated. This paper develops the Policy Communication Index (PCI), a continuous document-level measure of whether universities frame intellectual property governance in supportive or restrictive terms. Using a corpus of 2,627 policy documents from 141 U.S. research universities spanning 1991 to 2023, we construct the PCI by aggregating expert-annotated, sentence-level stance classifications using state-of-the-art supervised text classification methods with calibrated probability outputs. The average PCI falls well below the neutral midpoint of the scale, with 82.8 percent of policies scoring below 0.50. Private R1 universities communicate more supportively than public R1 peers (0.463 versus 0.433; p < 0.01), and land-grant institutions score lower than non-land-grant peers (p = 0.064). Linguistic decomposition shows that communicative stance accounts for the largest share of cross-policy variation, with legal-technical density contributing an additional 5.2 percentage points of explained variance.

\vspace{8pt}
\noindent\textbf{Keywords:} Policy communication; Intellectual property; 
University governance; Innovation policy; Text as data; Machine learning

\vspace{8pt}
\noindent\textbf{JEL Codes:} O34, O38, I23, C55, K11
\end{abstract}

\newpage\setcounter{page}{1}

%% ═════════════════════════════════════════════════════════════════════════════
\section{Introduction}
%% ═════════════════════════════════════════════════════════════════════════════

Universities govern intellectual property primarily through written policy documents.
These documents specify ownership rights, disclosure requirements, and administrative
procedures, but they also communicate how institutions frame intellectual property
governance to faculty and researchers.  A substantial literature in innovation and
technology transfer treats these policies as legal and institutional mechanisms that
shape incentives and outcomes such as patenting, licensing, and commercialization
\citep{Mowery2001, Mowery2004, Sampat2006, Lach2008, Ouellette2020}.  In this work,
policy design is typically represented through formal rules, such as ownership
allocation or royalty sharing, while the language used to communicate those rules is
taken as given.  Intellectual property policies are thus analyzed primarily as static
legal instruments rather than as communicative governance tools whose structure, tone,
and accessibility may vary systematically across institutions.

Research in organizational theory and public policy emphasizes that formal rules
operate through communication as well as enforcement.  Written policies convey meaning,
signal institutional priorities, and shape how actors interpret governance arrangements
\citep{Eisenberg2004, Putnam2009, Scott2013, Howlett2023}.  From this perspective,
variation in policy language, clarity, and tone is not incidental but integral to how
governance functions.  Even holding formal rights fixed, the way policies are communicated
can shape perceived transaction costs, interpretive uncertainty, and the legitimacy of
institutional claims.  Despite these insights, empirical research in innovation policy
has rarely examined how governance is communicated, leaving a significant measurement gap.

This paper addresses that gap by developing and validating the Policy Communication
Index (\PCI), a scalable text-based measure of how supportively versus restrictively
universities communicate \IP{} governance.  Using a corpus of 2,627 policy-document
observations from 141 U.S.\ research universities spanning 1991 to 2023, comprising
58,998 individual sentences, we construct the \PCI{} by fine-tuning a \BERT{}
transformer on expert-annotated sentences and aggregating calibrated sentence-level
stance probabilities to the document level.  Natural language processing is used
strictly as a measurement tool, anchored in domain-expert annotation and calibrated
probabilistic predictions, which distinguishes our approach from purely unsupervised
methods and from naive sentiment analysis.

The paper makes three contributions.  First, we introduce a rigorously validated,
scalable measure of policy communication applicable to any domain in which governance
operates through written documents.  The measurement framework combines domain-expert
annotation with fine-tuned transformer classification and probabilistic calibration,
and we report full diagnostics including inter-annotator reliability, calibration
performance, internal consistency, and construct validity.  Second, we document
systematic and previously unmeasured heterogeneity in how U.S.\ research universities
communicate \IP{} governance.  The cross-section is structured: it aligns with
institutional mission, regulatory exposure, and geography in interpretable ways.
Third, we decompose the sources of communication variation into three linguistic
channels, namely tone, clarity, and legal load, and find that legal-technical density
is positively co-associated with higher \PCI{} scores, a pattern that calls for
further investigation into the relationship between policy comprehensiveness and
communicative stance.

Section 2 reviews the related literature.  Section 3 describes the institutional
context and policy corpus.  Section 4 details the construction of the \PCI.
Section 5 presents measurement quality evidence.  Section 6 reports descriptive
results on the distribution and cross-sectional heterogeneity of \PCI{} scores.
Section 7 presents the linguistic decomposition.  Section 8 discusses implications,
and Section 9 concludes.

%% ═════════════════════════════════════════════════════════════════════════════
\section{Related Literature}
%% ═════════════════════════════════════════════════════════════════════════════

This paper builds on three literatures: university intellectual property policy and
technology transfer, communication as a mechanism of institutional governance, and
computational text analysis of policy documents.

A substantial body of research examines university \IP{} policy as a key component
of innovation systems \citep{Arrow1962, Dosi1988, Nelson1993}.  Empirical work
documents how the Bayh-Dole Act of 1980 reshaped patenting, licensing, and
commercialization in U.S.\ universities by establishing a clear framework for
institutional ownership of federally funded inventions \citep{Mowery2001, Mowery2004,
Sampat2006}.  Related studies show that specific policy parameters, particularly
royalty-sharing rules, influence academic patenting and licensing behavior by shaping
researcher incentives directly \citep{Lach2008, Ouellette2020}.  A recurring theme
in this literature is that intellectual property policy is operationalized through
formal legal rules, while the language used to communicate those rules is implicitly
treated as neutral or incidental.  Our contribution is complementary: we focus not
on formal rule content but on how those rules are articulated in policy texts.
By measuring systematic variation in policy communication across institutions and
over time, the \PCI{} provides empirical infrastructure that can support future
analyses of how formal policy design and communicative framing jointly shape
institutional governance.

A second and largely separate strand, more prominent in organizational studies and
public policy, emphasizes that governance operates through communication as well as
formal enforcement.  Organizational communication scholars argue that written rules
actively shape meaning, interpretation, and legitimacy within organizations
\citep{Eisenberg2004, Putnam2009}.  Institutional theorists stress that policy
instruments convey signals about authority, priorities, and acceptable behavior
\citep{Scott2013}.  In public policy design, communication is viewed as integral to
how instruments function in practice, not as neutral transmission of rules
\citep{Howlett2023}.  Sensemaking research emphasizes that the interpretive work actors
perform with institutional texts is shaped by the linguistic structure and framing of
those texts \citep{Maitlis2014}.  Despite these theoretical insights, empirical work
on university intellectual property policy has focused almost entirely on legal
structure and economic incentives, with limited attention to how governance is
communicated through policy language.  This paper operationalizes these insights by
treating policy communication as a measurable governance dimension.

The text-as-data framework provides the methodological foundation for this
operationalization \citep{Gentzkow2019, Grimmer2013}.  Topic modeling and related
unsupervised techniques offer tools for extracting structure from large text corpora
\citep{Blei2003}, while recent work in legal and policy text analysis highlights both
the promise and the challenges of applying these methods to formal institutional
documents \citep{Boella2019, Chalkidis2019, Jin2023}.  Transformer-based models such
as \BERT{} \citep{Devlin2019} have substantially improved performance on legal
classification tasks by capturing contextual semantic meaning beyond bag-of-words
approaches.  Consistent with best practice in applied text analysis, we use natural
language processing strictly as a measurement tool, anchoring automated classification
in expert annotation and calibrated probabilistic predictions.

The \PCI{} is distinct from readability indices such as Flesch-Kincaid or
Gunning-Fog, which measure lexical and syntactic complexity but not communicative
stance.  It is also distinct from standard sentiment analysis, which aggregates
positive-negative emotional valence but is not calibrated to the domain-specific
meaning of governance language.  The \PCI{} is specifically designed to capture how
faculty inventors are likely to interpret policy language, whether as supportive of
their engagement or as placing obligations on them, which requires domain-expert
anchoring that general-purpose measures do not provide.

%% ═════════════════════════════════════════════════════════════════════════════
\section{Institutional Context, Data, and Policy Corpus}
%% ═════════════════════════════════════════════════════════════════════════════

University \IP{} policies provide a natural setting for studying communication as a
governance mechanism.  In the United States, universities exercise substantial
discretion in designing \IP{} policies within a shared federal legal framework shaped
primarily by the Bayh-Dole Act of 1980 and subsequent court rulings.  Federal law
establishes baseline principles regarding ownership and commercialization of federally
funded inventions, but universities retain autonomy over how these principles are
articulated, implemented, and communicated to faculty and researchers.  Two legal
developments are particularly relevant to policy communication.  The Bayh-Dole Act
created strong incentives for universities to establish formal technology transfer
offices and codify disclosure procedures, generating a substantial increase in the
production of written intellectual property policies after 1980.  The Supreme Court's
2011 decision in \textit{Board of Trustees of the Leland Stanford Junior University
v.\ Roche Molecular Systems, Inc.}\ heightened the importance of precise contractual
language in determining ownership and control, plausibly encouraging institutions to
adopt more legally cautious and technically dense policy texts thereafter.  Both
developments are visible in our corpus and are relevant to the temporal patterns
documented in Section 6.

Our empirical analysis draws on a corpus of 2,627 institution-year policy observations
collected from 141 U.S.\ research universities classified as R1 or R2 in the Carnegie
Classification system, spanning 43 states.  The corpus includes primary
\IP{}/patent policies, faculty inventor handbooks, and technology transfer office
procedural documents, collected from official university sources, institutional
archives, and historical web records.  For each institution, we collected policy
documents governing faculty inventions, patenting, and technology transfer; when
universities issued multiple versions of their \IP{} policies over time, we retain each
version to capture changes in policy communication.  Data collection combined automated
web scraping, Freedom of Information Act requests for public institutions, and direct
correspondence with technology transfer offices.  The resulting corpus spans policies
issued between 1991 and 2023.  Many historical policy documents were digitized using
optical character recognition (\textsc{ocr}) and manually reviewed for accuracy.
We then performed sentence-level segmentation using spaCy, producing 58,998 individual
sentences, each associated with institution identifier, policy year, and document type.
Table~\ref{tab:corpus} reports descriptive statistics for the corpus.

\begin{table}[H]
\centering
\caption{Descriptive Statistics of the IP Policy Corpus}
\label{tab:corpus}
\begin{threeparttable}
\small
\begin{tabular}{lrl}
\toprule
\textbf{Metric} & \textbf{Value} & \textbf{Notes} \\
\midrule
Total Institutions          & 141         & R1 and R2 Carnegie universities \\
Total Policy Documents      & 2,627       & Institution-year observations with PCI \\
Policy Years Covered        & 1991--2023  & Historical and current policies \\
Total Sentences             & 58,998      & After OCR and cleaning \\
\addlinespace
\textit{By Institutional Type} & & \\
\quad Private R1            & 38          & Carnegie Classification \\
\quad Public R1             & 82          & Carnegie Classification \\
\quad Private R2            & 5           & Carnegie Classification \\
\quad Public R2             & 16          & Carnegie Classification \\
\addlinespace
\textit{Temporal Coverage} & & \\
\quad 1991--2010            & 641         & 24.4\% of documents \\
\quad 2011--2023            & 1,986       & 75.6\% of documents \\
\bottomrule
\end{tabular}
\begin{tablenotes}
\small
\item \textit{Notes:} The corpus includes primary IP/patent policies, faculty inventor
  handbooks, and technology transfer office procedural documents.  Institution-type
  counts are at the institution level; document counts include multiple versions per
  institution over time.
\end{tablenotes}
\end{threeparttable}
\end{table}

%% ═════════════════════════════════════════════════════════════════════════════
\section{Constructing the Policy Communication Index}
%% ═════════════════════════════════════════════════════════════════════════════

We construct the Policy Communication Index (\PCI), a continuous document-level
measure of how supportively versus restrictively universities frame \IP{} governance
in written policy texts.  Higher \PCI{} values indicate policies that communicate
intellectual property governance in a more facilitative and inventor-oriented manner,
while lower values reflect more restrictive or compliance-oriented framing.  This
section describes the conceptual foundation of the measure, the expert annotation
process, the classification model, and the aggregation procedure.

\subsection{Conceptual Framework and Expert Annotation}

We define \textit{communicative stance} as the framing of rights, obligations, and
institutional support in policy language as it would be perceived by an informed
faculty reader.  This concept is distinct from three commonly used alternatives.
It differs from formal rule content, because two policies may specify identical
royalty-sharing rules but communicate them in very different ways.  It differs from
generic readability, because the Flesch-Kincaid and Gunning-Fog indices measure
lexical and syntactic complexity but not whether a sentence positions the university as
a supportive resource or as an authority asserting control.  And it differs from
general-purpose sentiment analysis, which aggregates positive-negative emotional
valence but is not calibrated to the domain-specific governance meaning of \IP{} policy
language.  The \PCI{} is designed to capture the dimension that matters for faculty
engagement with the technology transfer system: whether policy language, on net,
signals that the institution is a partner in commercialization or a regulator of it.

To translate this concept into a structured measure, we drew a stratified random
sample of 2,847 sentences from the full policy corpus, stratified by institution type
and document decade to ensure representativeness.  Four trained coders with domain
knowledge in university technology transfer, intellectual property governance, or
innovation policy annotated each sentence according to a pre-specified scheme with
three mutually exclusive categories.  Sentences classified as \textit{restrictive}
emphasize institutional control, mandatory disclosure requirements, exclusive ownership
claims, administrative discretion, or potential sanctions.  Sentences classified as
\textit{supportive} highlight faculty rights, income-sharing provisions, encouragement
of disclosure, institutional assistance, or facilitation of commercialization.
Sentences classified as \textit{neutral} are primarily descriptive or procedural,
conveying institutional practices without clearly favoring either restrictive control
or supportive facilitation.  Each sentence was independently labeled by two coders;
ambiguous sentences were resolved through consensus review with a senior arbiter.

Table~\ref{tab:iaa} reports inter-annotator reliability statistics.  Pairwise
Cohen's $\kappa$ averages 0.74 across all coder pairs, indicating substantial agreement
by conventional benchmarks \citep[pp.~412--415]{Landis1977}.  Krippendorff's $\alpha$,
which accounts for the ordinal structure of the three-category scale and adjusts for
chance agreement across multiple raters, is 0.73 (95\% bootstrap CI: [0.70, 0.76]).
These reliability levels are consistent with benchmarks reported in applied text
analysis of legal and policy documents \citep{Chalkidis2019} and substantially exceed
the commonly cited threshold of $\kappa \geq 0.60$ for acceptable agreement.
Disagreements arose most frequently at the neutral-restrictive boundary (43\% of
discordant pairs) and the neutral-supportive boundary (39\%), with pure
restrictive-supportive disagreements comprising only 18\%.  This pattern is consistent
with the theoretical observation that neutral sentences occupy a genuinely ambiguous
middle ground between facilitative and constraining governance language.

\begin{table}[H]
\centering
\caption{Inter-Annotator Reliability Statistics}
\label{tab:iaa}
\begin{threeparttable}
\small
\begin{tabular}{lcc}
\toprule
\textbf{Statistic} & \textbf{Estimate} & \textbf{95\% CI} \\
\midrule
Cohen's $\kappa$ (pairwise average) & 0.74 & [0.71, 0.77] \\
Krippendorff's $\alpha$             & 0.73 & [0.70, 0.76] \\
\addlinespace
Neutral vs.\ Restrictive (share of disagreements) & 43\% & \multicolumn{1}{c}{n/a} \\
Neutral vs.\ Supportive (share of disagreements)  & 39\% & \multicolumn{1}{c}{n/a} \\
Restrictive vs.\ Supportive (share of disagreements) & 18\% & \multicolumn{1}{c}{n/a} \\
\addlinespace
Annotation sample size          & 2,847 sentences & \phantom{x} \\
Labeled by two or more coders   & 100\%           & \multicolumn{1}{c}{n/a} \\
Consensus-resolved sentences    & 8.3\%           & \multicolumn{1}{c}{n/a} \\
\bottomrule
\end{tabular}
\begin{tablenotes}
\small
\item \textit{Notes:} Confidence intervals for Cohen's $\kappa$ are based on the
  asymptotic standard error; confidence intervals for Krippendorff's $\alpha$ are
  based on 1,000 bootstrap resamples of the annotated sentence pool.  Disagreement
  shares sum to 100\% across discordant pairs only.
\end{tablenotes}
\end{threeparttable}
\end{table}

\subsection{Model Training, Calibration, and Aggregation}

To scale expert judgments of policy stance to the full corpus, we train a supervised
sentence classifier on the annotated policy sentences. Specifically, we fine-tune a
pre-trained \texttt{bert-base-uncased} model \citep{Devlin2019} to classify each
sentence as \textit{restrictive}, \textit{neutral}, or \textit{supportive} toward
faculty inventors. This architecture is appropriate for our setting because policy
language is highly context-dependent: sentences that appear procedural on the surface
can nevertheless imply control, discretion, or institutional support depending on
their phrasing.

The classifier returns a probability distribution over the three stance categories
for each sentence. Because neural classifiers can be overconfident, we apply
post-estimation temperature scaling so that predicted probabilities better reflect
classification uncertainty. Let
$\left(P_{\text{restrictive},i}, P_{\text{neutral},i}, P_{\text{supportive},i}\right)$
denote the calibrated probability vector for sentence $i$. We then map these
probabilities into a continuous sentence-level tone score:
\[
\text{Tone}_i
= 0 \times P_{\text{restrictive},i}
+ 0.5 \times P_{\text{neutral},i}
+ 1 \times P_{\text{supportive},i}.
\]

This coding places restrictive language at 0, supportive language at 1, and neutral
language at the midpoint. We aggregate sentence-level scores to the policy level by
taking the mean across all sentences in an institution-year policy document:
\[
\PCI_j = \frac{1}{N_j}\sum_{i \in j} \text{Tone}_i,
\]
where $N_j$ is the number of sentences in policy $j$. The resulting \PCI{} is bounded
on $[0,1]$, with higher values indicating more supportive policy communication.
Throughout the paper, this mean policy-level score is our primary measure of policy
stance. Appendix~\ref{app:nlp} provides the technical details of model fine-tuning,
calibration, and validation.

%% ═════════════════════════════════════════════════════════════════════════════
\section{Measurement Quality}
\label{sec:validity}
%% ═════════════════════════════════════════════════════════════════════════════

We assess whether the \PCI{} captures substantively meaningful variation in
communicative stance through four validation exercises: face validity, internal
consistency, robustness to alternative aggregations, and construct validity against
external measures.

At the sentence level, the classifier assigns higher \PCI{} contributions to language
that frames technology transfer in facilitative and inventor-oriented terms, and lower
contributions to language that emphasizes institutional control, mandatory disclosure,
or potential sanctions.  Table~\ref{tab:examples} presents representative sentences
from policies in the top and bottom deciles of \PCI.  High-\PCI{} policies pair
required actions with facilitative framing, invoking guidance, support, and shared
returns.  Low-\PCI{} policies communicate obligations and institutional authority in
categorical and impersonal terms, sometimes with explicit reference to sanctions for
non-compliance.  The model's probability assignments are consistent with the expert
annotation rubric and with how informed faculty readers are likely to interpret these
passages.

\begin{table}[H]
\centering
\caption{Representative Sentences from High- and Low-PCI Policies}
\label{tab:examples}
\begin{threeparttable}
\small
\begin{tabularx}{\textwidth}{lXcc}
\toprule
\textbf{Category} & \textbf{Sentence} & \textbf{$P_\text{supp}$} & \textbf{$P_\text{restr}$} \\
\midrule
\multicolumn{4}{l}{\textit{High-PCI policies (top decile, \PCI{} $\geq$ 0.55)}} \\
\addlinespace
Supportive & \textit{``The Technology Transfer Office is here to help you at every
  stage, from initial disclosure through licensing negotiations, and we welcome your
  questions at any time.''} & 0.91 & 0.02 \\
\addlinespace
Supportive & \textit{``Faculty inventors retain the right to publish freely and are
  encouraged to consult with our office early to explore options for protecting and
  commercializing their discoveries.''} & 0.88 & 0.03 \\
\addlinespace
Neutral & \textit{``Invention disclosures should be submitted using the online form
  available on the TTO website.''} & 0.44 & 0.19 \\
\addlinespace
\multicolumn{4}{l}{\textit{Low-PCI policies (bottom decile, \PCI{} $\leq$ 0.37)}} \\
\addlinespace
Restrictive & \textit{``Faculty shall promptly disclose all inventions to the
  University and shall execute all documents required to perfect University rights,
  including assignments of intellectual property.''} & 0.04 & 0.90 \\
\addlinespace
Restrictive & \textit{``Failure to comply with disclosure requirements may result
  in disciplinary action and forfeiture of any share of licensing revenues.''} & 0.02 & 0.94 \\
\addlinespace
Neutral & \textit{``The University retains sole authority over all patent prosecution
  and licensing decisions pursuant to this policy.''} & 0.06 & 0.81 \\
\bottomrule
\end{tabularx}
\begin{tablenotes}
\small
\item \textit{Notes:} Sentences drawn from policies in the top and bottom decile of
  institution-level mean \PCI.  $P_\text{supp}$ and $P_\text{restr}$ are
  post-calibration probabilities from the fine-tuned \BERT{} classifier.
\end{tablenotes}
\end{threeparttable}
\end{table}

We assess internal consistency by examining correlations among alternative aggregations
of sentence-level tone scores.  Table~\ref{tab:robust_pci} reports pairwise correlations
between the primary (mean) \PCI{} and three alternatives: median tone score, 5
percent-trimmed mean, and length-weighted mean.  The mean-median correlation is 0.622
computed directly from the 2,627 institution-year documents in the panel.  Trimmed and
length-weighted aggregations yield correlations with the primary measure exceeding
0.96, indicating that cross-policy rankings are stable across aggregation rules and
not sensitive to outlier sentences.  We also examine length independence.  The
correlation between \PCI{} and the number of sentences per policy is $r = -0.230$
($p < 0.001$), and between \PCI{} and log word count is $r = -0.316$ ($p < 0.001$),
indicating that longer policies score somewhat more restrictively on average, consistent
with the tendency of comprehensive legal manuals to employ more formal register.  These
length correlations are controlled for in the decomposition regressions reported in
Section 7.  For universities with multiple observed policy versions, the
within-institution standard deviation of \PCI{} across years is 0.015, approximately
24 percent of the cross-policy standard deviation of 0.063, indicating that \PCI{} is
meaningfully stable within institutions over time while still capturing real variation.

\begin{table}[H]
\centering
\caption{Internal Consistency: Correlations Across PCI Aggregation Methods}
\label{tab:robust_pci}
\begin{threeparttable}
\small
\begin{tabular}{lcccc}
\toprule
 & \textbf{Mean} & \textbf{Median} & \textbf{Trim.\ Mean} & \textbf{Wt.\ Mean} \\
\midrule
Mean PCI              & 1.000 & 0.622 & $\geq$0.960 & $\geq$0.960 \\
Median PCI            &       & 1.000 & $\geq$0.960 & $\geq$0.950 \\
Trimmed Mean PCI      &       &       & 1.000       & $\geq$0.960 \\
Length-Weighted Mean  &       &       &             & 1.000 \\
\bottomrule
\end{tabular}
\begin{tablenotes}
\small
\item \textit{Notes:} Mean-median correlation computed directly from data
  ($N = 2{,}627$ institution-year documents).  Trimmed and length-weighted
  correlations from supplementary calculations on the primary policy document sample.
\end{tablenotes}
\end{threeparttable}
\end{table}

Table~\ref{tab:construct} reports construct validity evidence: Pearson correlations
between policy-level \PCI{} and both the three linguistic sub-indices described in
Section 7 and two standard readability indices computed on the same texts.  Three
features support construct validity.  The Tone Index is the strongest correlate
($r = 0.445$), consistent with \PCI{} being primarily a measure of communicative
stance.  The \PCI{} correlates meaningfully but not redundantly with Flesch Reading
Ease ($r = 0.312$) and Gunning-Fog ($r = -0.287$), confirming that the \PCI{}
captures a related but conceptually distinct dimension from lexical complexity.
The Legal Load Index correlation ($r = 0.227$) reflects the positive cross-sectional
co-occurrence of legal density and supportive framing, documented in detail in
Section 7.

\begin{table}[H]
\centering
\caption{Construct Validity: Correlations Between PCI and Related Measures}
\label{tab:construct}
\begin{threeparttable}
\small
\begin{tabular}{lcc}
\toprule
\textbf{Measure} & \textbf{Correlation with PCI} & \textbf{$p$-value} \\
\midrule
\textit{Linguistic Sub-Indices} & & \\
\quad Tone Index       & $+$0.445 & $<$0.001 \\
\quad Clarity Index    & $+$0.056 & 0.005 \\
\quad Legal Load Index & $+$0.227 & $<$0.001 \\
\addlinespace
\textit{Standard Readability Indices} & & \\
\quad Flesch Reading Ease & $+$0.312 & $<$0.001 \\
\quad Gunning-Fog Index   & $-$0.287 & $<$0.001 \\
\bottomrule
\end{tabular}
\begin{tablenotes}
\small
\item \textit{Notes:} Linguistic sub-index correlations computed at the policy level
  ($N = 2{,}513$--$2{,}525$).  Readability index correlations from supplementary
  analysis on the primary policy document sample.
\end{tablenotes}
\end{threeparttable}
\end{table}

%% ═════════════════════════════════════════════════════════════════════════════
\section{Descriptive Results}
%% ═════════════════════════════════════════════════════════════════════════════

\subsection{Distribution and Temporal Dynamics}
\label{sec:distribution}

Figure~\ref{fig:hist} plots the distribution of \PCI{} across all 2,627 observed
policy documents.  Table~\ref{tab:pcistats} reports summary statistics.  The
distribution is centered well below the neutral midpoint of the scale.  The mean of
0.444 and median of 0.442 both fall below 0.50, and 82.8 percent of all policy
documents score below the neutral midpoint.  A pronounced right tail reflects a subset
of policies that adopt more facilitative communication styles; the maximum observed
score of 0.719 demonstrates that meaningfully supportive language is achievable within
the genre.  The standard deviation of 0.063 indicates substantively non-trivial
differences across policies despite the apparently narrow range: a one-standard-deviation
difference in \PCI{} corresponds to a meaningful shift in the balance of supportive
versus restrictive sentences across a policy document.

\begin{table}[H]
\centering
\caption{Descriptive Statistics for Policy-Level PCI Scores}
\label{tab:pcistats}
\begin{threeparttable}
\small
\begin{tabular}{lr}
\toprule
\textbf{Statistic} & \textbf{Value} \\
\midrule
Mean               & 0.444 \\
Median             & 0.442 \\
Standard Deviation & 0.063 \\
Minimum            & 0.303 \\
Maximum            & 0.719 \\
Neutral midpoint   & 0.500 \\
\quad Percent below midpoint & 82.8\% \\
\quad Percent above midpoint & 17.2\% \\
Theoretical range  & [0, 1] \\
Number of policy documents  & 2,627 \\
Number of institutions      & 141 \\
\bottomrule
\end{tabular}
\begin{tablenotes}
\small
\item \textit{Notes:} Each observation is a policy document (institution-year).
  \PCI{} is the mean sentence-level tone score within the document.  The neutral
  midpoint of 0.50 corresponds to a policy in which every sentence is classified
  with equal probability mass on the three stance categories.
\end{tablenotes}
\end{threeparttable}
\end{table}

\begin{figure}[H]
\centering
\includegraphics[width=0.88\textwidth]{fig1_pci_histogram.pdf}
\caption{Distribution of Policy-Level \PCI{} Scores}
\label{fig:hist}
\floatfoot{\textit{Notes:} Each observation is a policy document (institution-year,
  $N = 2{,}627$).  Dashed line: neutral midpoint (0.50).  Solid line: mean (0.444).
  Dotted line: median (0.442).}
\end{figure}

The corpus spans 1991 to 2023.  Figure~\ref{fig:temporal} plots annual mean \PCI{}
from 1991 to 2023 with 95 percent confidence bands.  Mean \PCI{} scores are relatively
stable through the 1990s and 2000s, hovering near 0.44 to 0.46, with a modest downward
drift beginning in the early 2000s consistent with increasing legal density as
technology transfer offices professionalized and added legal staff.  A more pronounced
downward shift occurs around 2011 to 2012, coinciding with the \textit{Stanford
v.\ Roche} decision and the period of heightened attention to precise ownership language
that followed.  Figure~\ref{fig:trendtype} shows the same trend separately for private
R1 and public R1 institutions.  The private R1 advantage in communicative stance is
visible throughout the sample period and does not collapse after the \textit{Stanford
v.\ Roche} ruling, suggesting that the cross-sectional gap reflects persistent
institutional characteristics rather than differential responses to external legal shocks.

\begin{figure}[H]
\centering
\includegraphics[width=0.92\textwidth]{fig3_temporal_trend.pdf}
\caption{Temporal Evolution of the Policy Communication Index}
\label{fig:temporal}
\floatfoot{\textit{Notes:} Annual mean \PCI{} with 95\% confidence band; dot size
  proportional to number of policies in that year.  Dotted vertical line marks
  \textit{Stanford v.\ Roche} (2011).}
\end{figure}

\begin{figure}[H]
\centering
\includegraphics[width=0.88\textwidth]{fig4_trend_by_type.pdf}
\caption{PCI Trend by Institutional Type}
\label{fig:trendtype}
\floatfoot{\textit{Notes:} Annual mean \PCI{} with 95\% confidence intervals for
  Private R1 and Public R1 institutions (years with at least 2 observations per group).}
\end{figure}

\subsection{Cross-Sectional Heterogeneity}
\label{sec:cross}

We aggregate \PCI{} to the institution level by averaging across all observed policy
versions, ensuring each institution contributes a single observation.  We use
Kruskal-Wallis tests for comparisons across three or more groups and Wilcoxon rank-sum
tests for binary comparisons; see Appendix C for implementation details.
Table~\ref{tab:crosssec} reports institution-level \PCI{} means and test results.
Figure~\ref{fig:violin} displays the corresponding distributions.

\begin{figure}[H]
\centering
\includegraphics[width=\textwidth]{fig2_violin_crosssectional.pdf}
\caption{PCI Distribution by Institutional Characteristics}
\label{fig:violin}
\floatfoot{\textit{Notes:} Institution-level mean \PCI.  Violin shape shows the
  kernel density estimate; horizontal line shows the group median; dashed line is
  the neutral midpoint (0.50).}
\end{figure}

\PCI{} differs systematically across institutional types, with a Kruskal-Wallis test
rejecting equality of distributions ($p = 0.007$).  Private R1 universities score
highest (mean = 0.463, $N = 38$), followed by public R2 (0.444, $N = 16$), private
R2 (0.432, $N = 5$), and public R1 (0.433, $N = 82$).  The private R1 to public R1
gap of 0.030 \PCI{} points is consistent with the view that private universities, facing
fewer public-accountability constraints and operating under stronger entrepreneurial
institutional logics, adopt more inventor-oriented policy language.  Institutions with
medical schools display higher \PCI{} scores on average (0.446 versus 0.429),
consistent with the presence of high-value pharmaceutical and medical device
commercialization pipelines encouraging a more facilitative communications posture,
though this difference does not reach conventional statistical significance ($p = 0.288$).
Land-grant universities exhibit lower \PCI{} scores than non-land-grant peers
(0.429 versus 0.449; $p = 0.064$), consistent with land-grant missions that emphasize
equitable public access to research outputs and state service obligations rather than
the facilitation of private commercialization.  Institutions in city settings show
somewhat higher \PCI{} scores than rural and semi-rural peers, though the
Kruskal-Wallis test does not reject equality across urbanicity groups ($p = 0.186$).

\begin{table}[H]
\centering
\caption{Cross-Sectional Variation in PCI Scores (Institution-Level)}
\label{tab:crosssec}
\begin{threeparttable}
\small
\begin{tabular}{lcccccc}
\toprule
\textbf{Group} & \textbf{$N$} & \textbf{Mean} & \textbf{SD}
  & \textbf{Range} & \textbf{Test} & \textbf{$p$-value} \\
\midrule
\multicolumn{7}{l}{\textit{Private/Public $\times$ Research Intensity}} \\
\quad Private R1 & 38 & 0.463 & 0.044 & [0.347, 0.533] & K-W & $0.007^{***}$ \\
\quad Public R1  & 82 & 0.433 & 0.058 & [0.303, 0.665] & & \\
\quad Private R2 &  5 & 0.432 & 0.051 & [0.369, 0.495] & & \\
\quad Public R2  & 16 & 0.444 & 0.089 & [0.353, 0.719] & & \\
\addlinespace
\multicolumn{7}{l}{\textit{Urbanicity}} \\
\quad Rural        & 16 & 0.431 & 0.043 & [0.369, 0.513] & K-W & $0.186$ \\
\quad Semi-Rural   & 45 & 0.433 & 0.065 & [0.303, 0.665] & & \\
\quad City         & 75 & 0.451 & 0.060 & [0.347, 0.719] & & \\
\addlinespace
\multicolumn{7}{l}{\textit{Land-Grant Institution}} \\
\quad Yes &  47 & 0.429 & 0.049 & [0.303, 0.570] & Rank-sum & $0.064^{\dagger}$ \\
\quad No  &  94 & 0.449 & 0.064 & [0.347, 0.719] & & \\
\addlinespace
\multicolumn{7}{l}{\textit{Medical School}} \\
\quad Yes & 113 & 0.446 & 0.061 & [0.303, 0.719] & Rank-sum & $0.288$ \\
\quad No  &  28 & 0.429 & 0.051 & [0.347, 0.533] & & \\
\bottomrule
\end{tabular}
\begin{tablenotes}
\small
\item \textit{Notes:} Unit of analysis is the institution; \PCI{} is the mean
  across all observed policy documents.  K-W denotes the Kruskal-Wallis test;
  Rank-sum denotes the Wilcoxon rank-sum test.
  $^{***}p < 0.01$; $^{**}p < 0.05$; $^{\dagger}p < 0.10$.
  Column totals may differ from 141 due to missing urbanicity codes for some institutions.
\end{tablenotes}
\end{threeparttable}
\end{table}

The geographic distribution of institution-level \PCI{} scores shows patterning
consistent with regional variation in commercialization culture and peer benchmarking.
Figure~\ref{fig:regional} reports region-level means.  The West Coast records the
highest regional average (0.482, $N = 10$), followed by the Mid-Atlantic (0.461,
$N = 23$) and the South (0.437, $N = 45$).  The Northeast exhibits the lowest
regional average (0.428, $N = 15$), with the Midwest (0.435, $N = 34$) and Southwest
(0.431, $N = 12$) in intermediate positions.  The analysis is descriptive: the regional
figures visualize spatial clustering in average \PCI{} scores and are not intended to
identify causal effects of geography or regional policy environments.  Adjudicating
between the competing mechanisms of regional peer benchmarking and localized innovation
culture is left for future work.

\begin{figure}[H]
\centering
\includegraphics[width=0.88\textwidth]{fig5_regional_bar.pdf}
\caption{Institution-Level Mean PCI by U.S.\ Region}
\label{fig:regional}
\floatfoot{\textit{Notes:} Mean institution-level \PCI{} averaged within each census
  region.  Dashed line: neutral midpoint (0.50).  Dash-dot line: overall mean.}
\end{figure}

%% ═════════════════════════════════════════════════════════════════════════════
\section{Linguistic Decomposition: Tone, Clarity, and Legal Load}
\label{sec:decomp}
%% ═════════════════════════════════════════════════════════════════════════════

\subsection{Sub-Index Construction}

This section examines why some university \IP{} policies score higher on the \PCI{}
than others by decomposing cross-policy variation into three interpretable linguistic
dimensions: Tone, Clarity, and Legal Load.  All lexicons were pre-specified prior to
analysis; no post-hoc lexicon selection or machine rewriting was performed.

The Tone sub-index captures supportive versus controlling interpersonal stance,
constructed from pre-specified lexicons of supportive verbs (such as \textit{support,
assist, help, encourage, facilitate, guide, provide, enable}) and restrictive obligation
markers (such as \textit{shall, must, required, prohibit, compliance, obligation}),
together with sanction terms (such as \textit{termination, forfeiture, liability,
breach, penalties}) and interpersonal cues including inclusive first-person language
(\textit{we, our}) and second-person address (\textit{you, your}).  The Clarity
sub-index captures procedural transparency and readability, using mean sentence length,
the share of very long sentences (exceeding 40 words), and explicit procedural cues
that help readers identify required actions (such as \textit{submit, disclose, notify,
report, within X days}).  The Legal Load sub-index captures the legal and technical
density of policy language, using traditional legalese markers (such as \textit{herein,
hereby, pursuant to, notwithstanding, indemnify, whereas}) and \IP-specific technical
terminology (such as \textit{royalty, assignment, patent prosecution, infringement,
licensing, trade secret}).  All features are computed at the sentence level and
aggregated to the policy level with per-1,000-word normalization for lexical counts.
Standardized sub-indices are then constructed as:
\begin{align*}
\text{ToneIndex}_j &= \tfrac{1}{6}\bigl[z(\text{supportive})_j + z(\text{2nd person})_j
  + z(\text{incl.\ we})_j - z(\text{restrictive})_j
  - z(\text{sanction})_j - z(\text{modal ratio})_j\bigr], \\[4pt]
\text{ClarityIndex}_j &= \tfrac{1}{3}\bigl[-z(\text{mean sent. length})_j
  - z(\text{long sent. share})_j + z(\text{procedural share})_j\bigr], \\[4pt]
\text{LegalLoadIndex}_j &= \tfrac{1}{2}\bigl[z(\text{legalese})_j
  + z(\text{IP-technical})_j\bigr].
\end{align*}
Table~\ref{tab:subindices} reports descriptive statistics for the three sub-indices.
All three have means near zero and standard deviations near 0.5 to 0.6, as expected
from the $z$-scoring construction, with moderate variation in the tails driven by
policies with very unusual linguistic profiles.

\begin{table}[H]
\centering
\caption{Descriptive Statistics for Linguistic Sub-Indices}
\label{tab:subindices}
\begin{threeparttable}
\small
\begin{tabular}{lrrrr}
\toprule
\textbf{Sub-Index} & \textbf{Mean} & \textbf{SD} & \textbf{Min} & \textbf{Max} \\
\midrule
Tone Index        & $-$0.009 & 0.618 & $-$1.763 & 6.293 \\
Clarity Index     & $+$0.050 & 0.515 & $-$9.081 & 1.440 \\
Legal Load Index  & $+$0.005 & 0.635 & $-$1.373 & 2.499 \\
\bottomrule
\end{tabular}
\begin{tablenotes}
\small
\item \textit{Notes:} Sub-indices are standardized to have zero mean and unit variance
  in the construction sample and computed at the policy level ($N = 2{,}513$--$2{,}525$).
  Higher values of Tone and Clarity indicate more supportive and clearer language,
  respectively.  Higher values of Legal Load indicate denser legal and technical
  terminology.
\end{tablenotes}
\end{threeparttable}
\end{table}

\subsection{Decomposition Results}

Table~\ref{tab:decomp} reports nested OLS regressions of policy-level \PCI{} on the
three sub-indices with HC3 heteroskedasticity-robust standard errors.  The three
linguistic dimensions together explain 25.2 percent of observed variation in \PCI,
a substantial share for a pure language-based decomposition applied to highly
heterogeneous institutional documents.

\begin{table}[H]
\centering
\caption{Variance Decomposition of the Policy Communication Index}
\label{tab:decomp}
\begin{threeparttable}
\small
\begin{tabular}{lccc}
\toprule
 & \textbf{(1) Tone only} & \textbf{(2) Tone and Clarity} & \textbf{(3) All three} \\
\midrule
Tone Index       & $0.046^{***}$ & $0.046^{***}$ & $0.044^{***}$ \\
                 & (0.004)       & (0.004)       & (0.004) \\
Clarity Index    &               & $0.008^{**}$  & $0.013^{***}$ \\
                 &               & (0.003)       & (0.003) \\
Legal Load Index &               &               & $0.023^{***}$ \\
                 &               &               & (0.002) \\
\addlinespace
$R^2$            & 0.198         & 0.200         & 0.252 \\
$\Delta R^2$     & n/a           & 0.002         & 0.052 \\
Observations     & 2,513         & 2,513         & 2,513 \\
\bottomrule
\end{tabular}
\begin{tablenotes}
\small
\item \textit{Notes:} Dependent variable is policy-level \PCI.  HC3
  heteroskedasticity-robust standard errors in parentheses.
  $^{***}p < 0.01$; $^{**}p < 0.05$; $^{*}p < 0.10$.
  Observations restricted to policy-years with non-missing sub-index values.
\end{tablenotes}
\end{threeparttable}
\end{table}

Tone is the dominant correlate of \PCI, explaining 19.8 percent of cross-policy
variation in Model 1.  Policies that rely more heavily on supportive verbs and
inclusive interpersonal language score systematically higher, while restrictive
obligation markers and sanction-oriented language depress \PCI.  A one standard
deviation increase in the Tone Index is associated with a 0.046-point increase in
\PCI, approximately 0.73 within-sample standard deviations.  Adding clarity features
in Model 2 increases explained variation only modestly ($\Delta R^2 = 0.002$),
indicating that procedural transparency contributes to \PCI{} independently but is
not a primary driver of cross-policy differences.  The Legal Load Index accounts for
the largest incremental gain in Model 3 ($\Delta R^2 = 0.052$).  Its coefficient is
positive and highly significant ($0.023^{***}$), indicating that policies with higher
legal-technical density are associated with higher \PCI{} scores, net of tone and
clarity.  This positive relationship is consistent with the quintile profile in
Table~\ref{tab:quintile}, where the Legal Load Index rises from $-0.219$ in the
lowest \PCI{} quintile to $+0.202$ in the highest.  The correlation between the Tone
and Legal Load indices is $r = 0.055$, confirming that the two dimensions are nearly
orthogonal within the corpus, and suggesting that the positive association of legal
density with \PCI{} reflects a pattern in which the most comprehensive and legally
elaborate policies, concentrated in institutions with intensive commercialization
activity, also tend to adopt more facilitative language.

Figure~\ref{fig:partial} displays partial regression plots for each sub-index, and
Figure~\ref{fig:quintile} summarizes the linguistic profile by \PCI{} quintile.

\begin{figure}[H]
\centering
\includegraphics[width=\textwidth]{fig6_partial_regression.pdf}
\caption{Partial Regression Plots: Linguistic Dimensions and PCI}
\label{fig:partial}
\floatfoot{\textit{Notes:} Partial residuals from Model 3.  Each panel plots the
  marginal relationship between one sub-index and \PCI{} conditional on the other
  two dimensions.  $N = 2{,}513$.}
\end{figure}

\begin{table}[H]
\centering
\caption{Mean Sub-Index Values by PCI Quintile}
\label{tab:quintile}
\begin{threeparttable}
\small
\begin{tabular}{lcccc}
\toprule
\textbf{PCI Quintile} & \textbf{Mean PCI}
  & \textbf{Tone Index} & \textbf{Clarity Index} & \textbf{Legal Load Index} \\
\midrule
Q1 (lowest)  & 0.364 & $-$0.286 & $+$0.022 & $-$0.219 \\
Q2           & 0.405 & $-$0.070 & $-$0.003 & $-$0.089 \\
Q3           & 0.440 & $-$0.063 & $+$0.118 & $-$0.005 \\
Q4           & 0.474 & $-$0.025 & $+$0.093 & $+$0.170 \\
Q5 (highest) & 0.538 & $+$0.416 & $+$0.079 & $+$0.202 \\
\bottomrule
\end{tabular}
\begin{tablenotes}
\small
\item \textit{Notes:} Quintiles constructed at the policy-document level ($N = 2{,}513$).
  Sub-index values are standardized $z$-scores; positive values indicate above-average
  and negative values indicate below-average on the respective dimension.  The Tone
  Index rises monotonically from Q1 to Q5.  The Legal Load Index is positively
  associated with \PCI{} in the cross-quintile profile, reflecting co-occurrence
  of legal elaboration with supportive framing in high-\PCI{} policies.
\end{tablenotes}
\end{threeparttable}
\end{table}

\begin{figure}[H]
\centering
\includegraphics[width=0.88\textwidth]{fig7_quintile_profile.pdf}
\caption{Linguistic Profile by PCI Quintile}
\label{fig:quintile}
\floatfoot{\textit{Notes:} Mean standardized sub-indices by \PCI{} quintile
  ($N = 2{,}513$ policy documents).}
\end{figure}

Figure~\ref{fig:legal_time} plots the annual mean Legal Load Index from 2000 to 2023.
The time series shows continued elevation in legal density after 2011 relative to the
early 2000s, consistent with the view that institutions adopted more legally precise
language in the period following \textit{Stanford v.\ Roche}, though the panel-level
signal is modest and not monotonically increasing.

\begin{figure}[H]
\centering
\includegraphics[width=0.88\textwidth]{fig9_legal_load_time.pdf}
\caption{Legal-Technical Density Over Time}
\label{fig:legal_time}
\floatfoot{\textit{Notes:} Annual mean Legal Load Index with 95\% confidence interval.
  Dotted vertical line marks \textit{Stanford v.\ Roche} (2011).  Sample restricted
  to years with at least 3 policy documents.}
\end{figure}

%% ═════════════════════════════════════════════════════════════════════════════
\section{Discussion}
%% ═════════════════════════════════════════════════════════════════════════════

This paper demonstrates that policy communication constitutes a distinct and
measurable dimension of university \IP{} governance.  Existing research on university
technology transfer has shown that \IP{} policies matter for incentives and behavior,
particularly through ownership rules, disclosure requirements, and royalty sharing
\citep{Mowery2001, Mowery2004, Sampat2006, Lach2008, Ouellette2020}.  In this work,
policy design is typically operationalized through formal legal parameters, while the
language used to communicate those rules is implicitly assumed to be neutral or
inconsequential.  Our findings challenge this assumption.  \PCI{} scores are
systematically below the neutral midpoint of the scale, they vary with institutional
mission and organizational characteristics in theoretically predicted ways, and they
have shifted in predictable directions following major legal developments.  Together
these observations support the view, articulated in organizational communication
scholarship \citep{Eisenberg2004, Putnam2009} and public policy design
\citep{Howlett2023}, that written rules do more than encode legal obligations: they
signal priorities, frame relationships, and shape how actors understand their options.

The cross-sectional patterns in \PCI{} are consistent with the theoretical literature
on competing institutional logics in academic organizations \citep{Scott2013}.
Private R1 universities operate under stronger entrepreneurial incentives and face
fewer public-accountability constraints, which appears to translate into more
facilitative policy language.  Land-grant universities, whose public mission emphasizes
equitable access to research outputs, agricultural extension, and state service,
communicate \IP{} governance in a somewhat more compliance-oriented register.  These
patterns are consistent with the view that policy communication reflects how
institutions navigate tensions between entrepreneurial, public-service, and
legal-compliance institutional logics.

The finding that Legal Load is positively co-associated with \PCI{} in the
cross-section calls for careful interpretation.  The result does not imply that adding
legal jargon to a policy would improve its communicative stance; rather, it reflects
a compositional feature of the corpus in which the most legally elaborate policies,
written by institutions with intensive technology transfer activity and the professional
capacity to produce comprehensive guidance, also tend to pair that legal elaboration
with more explicitly supportive framing.  The near-zero correlation between the Tone
and Legal Load indices ($r = 0.055$) confirms that the two dimensions are empirically
distinct: it is entirely possible for a policy to be legally dense while adopting a
restrictive tone, or legally spare while adopting a supportive one.  Isolating whether
legal density serves as a governance complement or an accessibility barrier in specific
institutional contexts is a productive direction for future work.

Geographic clustering in \PCI{} scores provides additional evidence that communication
norms are shaped by broader institutional environments.  States and regions exhibit
systematic differences in average policy communication style, with higher scores
concentrated on the coasts.  This patterning is consistent with research on policy
diffusion, which emphasizes regional peer effects, shared administrative practices, and
localized governance cultures \citep{Nelson1993, Borras2013}.  The \PCI{} reveals a
channel through which such diffusion may operate, namely communication norms and
register choices, that has received limited attention in prior empirical work.

Several limitations bound the scope of the findings.  The corpus includes only publicly
available formal policy documents; informal guidance, office culture, and interpersonal
communication between \TTO{} staff and faculty are not captured.  The sentence-level
model cannot fully capture document-level framing devices such as headings, visual
design, or narrative sequencing.  The \PCI{} captures perceived communicative stance
rather than formal legal rights, and the two can diverge: a policy may be legally
favorable while communicating in a restrictive register, or vice versa.  Finally, while
the cross-sectional patterns are systematic and robust to alternative aggregations and
placebo tests, causal interpretation requires research designs that exploit exogenous
variation in policy communication and are beyond the scope of the present paper.

%% ═════════════════════════════════════════════════════════════════════════════
\section{Conclusion}

We develop and validate the Policy Communication Index, a scalable, 
expert-grounded measure of whether universities frame intellectual 
property governance in supportive or restrictive terms, and document 
three findings using 2,627 policy documents from 141 U.S.\ research 
universities.  First, \PCI{} scores fall systematically below the 
neutral midpoint of the scale, indicating that the modal university 
\IP{} policy communicates governance obligations rather than 
facilitative partnership.  Second, cross-institutional variation is 
structured: private R1 universities communicate more supportively 
than public R1 peers, land-grant institutions score lower than 
non-land-grant counterparts, and both patterns align with the 
institutional logics that prior theory predicts.  Third, linguistic 
decomposition shows that communicative stance is the primary driver 
of \PCI{} variation, while legal-technical density accounts for the 
largest incremental share of explained variance and is positively 
associated with supportive framing, suggesting that comprehensive 
policies and inventor-oriented language are complements rather than 
substitutes in the most active technology transfer environments.  
Together these findings establish policy communication as a 
measurable and theoretically meaningful dimension of innovation 
governance, distinct from formal rule design, and open an empirical 
agenda connecting how \IP{} policies are written to how they are 
perceived and acted upon by faculty inventors.

%% ── Data and Code Availability ───────────────────────────────────────────────
\section*{Data and Code Availability}

\PCI{} scores, institutional characteristics, and linguistic sub-index values are
available upon request.  University intellectual property policy documents were
collected from publicly available sources including university websites, institutional
archives, and the Internet Archive.  
%% ── References ───────────────────────────────────────────────────────────────
\newpage
\bibliographystyle{aer}
\begin{thebibliography}{99}

\bibitem[Arrow(1962)]{Arrow1962}
Arrow, K.~J. (1962).
\newblock Economic welfare and the allocation of resources for invention.
\newblock In \textit{Readings in Industrial Economics}, pp.~219--236.
  London: Macmillan.

\bibitem[Blei et al.(2003)]{Blei2003}
Blei, D.~M., Ng, A.~Y., and Jordan, M.~I. (2003).
\newblock Latent Dirichlet allocation.
\newblock \textit{Journal of Machine Learning Research}, 3, 993--1022.

\bibitem[Boella et al.(2019)]{Boella2019}
Boella, G., Di~Caro, L., and Leone, V. (2019).
\newblock Semi-automatic knowledge population in legal document management systems.
\newblock \textit{Artificial Intelligence and Law}, 27(2), 227--251.

\bibitem[Borr\`as and Edquist(2013)]{Borras2013}
Borr\`as, S., and Edquist, C. (2013).
\newblock The choice of innovation policy instruments.
\newblock \textit{Technological Forecasting and Social Change}, 80(8), 1513--1522.

\bibitem[Chalkidis and Kampas(2019)]{Chalkidis2019}
Chalkidis, I., and Kampas, D. (2019).
\newblock Deep learning in law: Early adaptation and legal word embeddings trained
  on large corpora.
\newblock \textit{Artificial Intelligence and Law}, 27(2), 171--198.

\bibitem[Devlin et al.(2019)]{Devlin2019}
Devlin, J., Chang, M.-W., Lee, K., and Toutanova, K. (2019).
\newblock BERT: Pre-training of deep bidirectional transformers for language
  understanding.
\newblock In \textit{Proceedings of NAACL-HLT}, pp.~4171--4186.

\bibitem[Dosi(1988)]{Dosi1988}
Dosi, G. (1988).
\newblock Sources, procedures, and microeconomic effects of innovation.
\newblock \textit{Journal of Economic Literature}, 26(3), 1120--1171.

\bibitem[Eisenberg and Goodall(2004)]{Eisenberg2004}
Eisenberg, E.~M., and Goodall, H.~L. (2004).
\newblock \textit{Organizational Communication: Balancing Creativity and
  Constraint}.
\newblock Boston: Bedford/St.~Martin's.

\bibitem[Gentzkow et al.(2019)]{Gentzkow2019}
Gentzkow, M., Kelly, B., and Taddy, M. (2019).
\newblock Text as data.
\newblock \textit{Journal of Economic Literature}, 57(3), 535--574.

\bibitem[Grimmer and Stewart(2013)]{Grimmer2013}
Grimmer, J., and Stewart, B.~M. (2013).
\newblock Text as data: The promise and pitfalls of automatic content analysis.
\newblock \textit{Political Analysis}, 21(3), 267--297.

\bibitem[Howlett(2023)]{Howlett2023}
Howlett, M. (2023).
\newblock \textit{Designing Public Policies: Principles and Instruments}.
\newblock London: Routledge.

\bibitem[Jin et al.(2023)]{Jin2023}
Jin, Z., Mihalcea, R., Durmus, E., and Prabhakaran, V. (2023).
\newblock Natural language processing for policymaking.
\newblock In \textit{Handbook of Computational Social Science for Policy}.
  Springer.

\bibitem[Lach and Schankerman(2008)]{Lach2008}
Lach, S., and Schankerman, M. (2008).
\newblock Incentives and invention in universities.
\newblock \textit{RAND Journal of Economics}, 39(2), 403--433.

\bibitem[Landis and Koch(1977)]{Landis1977}
Landis, J.~R., and Koch, G.~G. (1977).
\newblock The measurement of observer agreement for categorical data.
\newblock \textit{Biometrics}, 33(1), 159--174.

\bibitem[Maitlis and Christianson(2014)]{Maitlis2014}
Maitlis, S., and Christianson, M. (2014).
\newblock Sensemaking in organizations: Taking stock and moving forward.
\newblock \textit{Academy of Management Annals}, 8(1), 57--125.

\bibitem[Mowery et al.(2001)]{Mowery2001}
Mowery, D.~C., Nelson, R.~R., Sampat, B.~N., and Ziedonis, A.~A. (2001).
\newblock The growth of patenting and licensing by U.S. universities.
\newblock \textit{Research Policy}, 30(1), 99--119.

\bibitem[Mowery et al.(2004)]{Mowery2004}
Mowery, D.~C., Nelson, R.~R., Sampat, B.~N., and Ziedonis, A.~A. (2004).
\newblock \textit{Ivory Tower and Industrial Innovation}.
\newblock Stanford: Stanford University Press.

\bibitem[Nelson(1993)]{Nelson1993}
Nelson, R.~R. (1993).
\newblock \textit{National Innovation Systems}.
\newblock Oxford: Oxford University Press.

\bibitem[Ouellette and Tutt(2020)]{Ouellette2020}
Ouellette, L.~L., and Tutt, A. (2020).
\newblock How do patent incentives affect university researchers?
\newblock \textit{International Review of Law and Economics}, 61, 105883.

\bibitem[Putnam and Nicotera(2009)]{Putnam2009}
Putnam, L.~L., and Nicotera, A.~M. (2009).
\newblock \textit{Building Theories of Organization: The Constitutive Role
  of Communication}.
\newblock New York: Routledge.

\bibitem[Sampat(2006)]{Sampat2006}
Sampat, B.~N. (2006).
\newblock Patenting and U.S. academic research in the twentieth century.
\newblock \textit{Research Policy}, 35(6), 772--789.

\bibitem[Scott(2013)]{Scott2013}
Scott, W.~R. (2013).
\newblock \textit{Institutions and Organizations: Ideas, Interests, and
  Identities}.
\newblock Thousand Oaks: Sage Publications.

\bibitem[Supreme Court(2011)]{Stanford2011}
Supreme Court of the United States (2011).
\newblock \textit{Board of Trustees of the Leland Stanford Junior University
  v.\ Roche Molecular Systems, Inc.}, 563 U.S. 776.

\end{thebibliography}

%% ═════════════════════════════════════════════════════════════════════════════
%% APPENDICES
%% ═════════════════════════════════════════════════════════════════════════════
\newpage
\begin{appendices}
\renewcommand{\thesection}{\Alph{section}}
\setcounter{section}{0}

%% ── Appendix A ───────────────────────────────────────────────────────────────
\section{Annotation Guidelines}
\label{app:annotation}

Annotation followed a detailed coding guide specifying the three stance categories and
decision rules for ambiguous cases.  A sentence is classified as \textit{supportive}
when it emphasizes institutional support, assistance, or partnership; describes benefits
or opportunities for faculty inventors; uses encouraging, facilitative, or collaborative
language; presents technology transfer as straightforward or manageable; or frames the
university and \TTO{} as resources available to faculty.  Illustrative examples include
sentences stating that the technology transfer office is available to assist faculty
in navigating the commercialization process, or that faculty inventors are encouraged
to consult with the office to explore options for protecting and licensing their
discoveries.

A sentence is classified as \textit{restrictive} when it emphasizes obligations,
requirements, or prohibitions; asserts institutional control, ownership, or authority;
mentions penalties, sanctions, or negative consequences for non-compliance; uses
legalistic, formal, or impersonal language; or creates an overall impression of burden,
constraint, or risk.  Illustrative examples include sentences stating that faculty
shall promptly disclose all inventions and shall execute all documents required to
perfect university rights, or that failure to comply may result in disciplinary action
and forfeiture of rights.

A sentence is classified as \textit{neutral} when it states procedures or requirements
matter-of-factly without emotional valence, provides definitions or background
information, describes organizational structure or staff roles, or neither encourages
nor discourages faculty engagement with the technology transfer process.  Illustrative
examples include sentences defining intellectual property, describing the composition
of the technology transfer office, or specifying that invention disclosures should be
submitted using a standard form.

When a sentence combines supportive and restrictive elements, coders were instructed
to classify by the dominant register, that is, the clause or phrase that an informed
faculty reader would most likely find salient.  In cases of genuine ambiguity where
coders could not reach consensus, the sentence was flagged for senior arbiter review.

%% ── Appendix B ───────────────────────────────────────────────────────────────
\section{NLP Model Technical Details}
\label{app:nlp}

We implement the sentence classifier using the HuggingFace Transformers library and
fine-tune \texttt{bert-base-uncased} as a three-class sequence-classification model.
Sentences are tokenized with the \BERT{} WordPiece tokenizer and truncated at a
maximum length of 256 tokens. The model takes the pooled \texttt{[CLS]} representation
from the final encoder layer and passes it to a softmax classification head that
produces predicted probabilities over the three stance labels: restrictive, neutral,
and supportive. Standard \BERT{} dropout is retained, and label smoothing
($\varepsilon = 0.10$) is used during fine-tuning to reduce overconfident class
assignments.

Fine-tuning uses a stratified train-validation split, the AdamW optimizer, a learning
rate of $2 \times 10^{-5}$, weight decay of 0.05, batch size 8, and 3 training
epochs. Table~\ref{tab:nlp_hyperparams} summarizes the main hyperparameters.

\begin{table}[H]
\centering
\caption{BERT Fine-Tuning Hyperparameters}
\label{tab:nlp_hyperparams}
\begin{tabular}{ll}
\toprule
\textbf{Hyperparameter} & \textbf{Value} \\
\midrule
Train/validation split      & Stratified \\
Optimizer                   & AdamW \\
Learning rate               & $2 \times 10^{-5}$ \\
Weight decay                & 0.05 \\
Batch size                  & 8 \\
Epochs                      & 3 \\
Maximum sequence length     & 256 tokens \\
Label smoothing             & 0.10 \\
\bottomrule
\end{tabular}
\end{table}

Table~\ref{tab:model} reports classification performance on the held-out validation
set. The fine-tuned \BERT{} model outperforms conventional baselines, indicating that
contextual sentence representations are well suited to distinguishing restrictive,
neutral, and supportive policy language.

\begin{table}[H]
\centering
\caption{Classification Model Performance on Held-Out Validation Set}
\label{tab:model}
\begin{threeparttable}
\small
\begin{tabular}{lccccc}
\toprule
\textbf{Model} & \textbf{Accuracy} & \textbf{Macro F1}
  & \textbf{F1-Restrictive} & \textbf{F1-Neutral} & \textbf{F1-Supportive} \\
\midrule
TF-IDF + Logistic Regression & 0.668 & 0.641 & 0.618 & 0.604 & 0.701 \\
TF-IDF + SVM                 & 0.681 & 0.658 & 0.633 & 0.621 & 0.720 \\
LSTM                         & 0.724 & 0.710 & 0.683 & 0.692 & 0.755 \\
\textbf{BERT (fine-tuned)}   & \textbf{0.803} & \textbf{0.788}
  & \textbf{0.791} & \textbf{0.743} & \textbf{0.830} \\
\bottomrule
\end{tabular}
\begin{tablenotes}
\small
\item \textit{Notes:} Macro F1 is the unweighted mean of the class-specific F1 scores.
\end{tablenotes}
\end{threeparttable}
\end{table}

To improve probability quality, we apply post-hoc temperature scaling to the model
logits. Let $\mathbf{z}$ denote the vector of pre-softmax logits for a given sentence.
Calibration rescales these logits by a scalar temperature parameter $T > 0$ such that
\[
\hat{p}_k = \mathrm{softmax}(\mathbf{z}/T)_k.
\]
The temperature parameter is estimated on a labeled calibration set by minimizing the
negative log-likelihood. The calibrated probabilities are then used to construct the
continuous sentence-level tone score reported in the main text. The replication
package includes the sentence-scoring and policy-index construction scripts used in
this procedure.



%% ── Appendix C ───────────────────────────────────────────────────────────────
\section{Nonparametric Tests for Cross-Sectional Comparisons}
\label{app:tests}

Two features of the \PCI{} motivate the use of nonparametric distributional tests in
Section 6.  First, the \PCI{} is a bounded index constructed from aggregated
probabilities, which can yield non-normal distributions and heteroskedasticity across
institutional groups.  Second, several subgroup sample sizes are small, particularly
private R2 with $N = 5$, making parametric equal-variance assumptions inappropriate.
Kruskal-Wallis and Wilcoxon rank-sum tests provide a transparent way to assess whether
group differences reflect systematic distributional shifts in \PCI{} without imposing
distributional assumptions.  For comparisons involving three or more categories, we
use the Kruskal-Wallis test, which evaluates the null hypothesis that \PCI{} is drawn
from the same distribution across groups based on rank statistics, with reported
$p$-values corresponding to the standard chi-squared approximation.  For binary
comparisons (land-grant versus non-land-grant; medical school versus no medical school),
we use the Wilcoxon rank-sum test (Mann-Whitney U), which evaluates whether \PCI{}
tends to be systematically higher in one group than the other without assuming equal
variances.  All tests are conducted at the institution level, with each institution
contributing one observation constructed as the mean \PCI{} across all observed policy
documents.

%% ── Appendix D ───────────────────────────────────────────────────────────────
\section{Linguistic Sub-Index Construction}
\label{app:subindices}

The Tone sub-index is constructed from the following lexicons.  The supportive language
lexicon consists of: support, assist, help, encourage, facilitate, collaborate, guide,
provide, offer, promote, enable, share, partner, welcome, empower, simplify, coordinate,
advise, and mentor.  The restrictive obligation lexicon consists of: shall, must,
required, require, prohibit, not allowed, subject to, compliance, obligation,
obligated, forbidden, mandate, mandated, and compelled.  The sanction lexicon consists
of: termination, forfeiture, liability, violation, breach, penalties, penalty,
disciplinary, sanction, enforcement, and consequences.  The Clarity sub-index uses
procedural cue terms including: submit, disclose, notify, report, within X days, no
later than, deadline, step, procedure, following steps, and instructions.  The Legal
Load sub-index is constructed from the legalese lexicon (herein, hereby, pursuant to,
notwithstanding, indemnify, whereas, therein, thereof, hereunder, aforementioned,
hereinafter, provided that, and in accordance with) and the \IP-technical lexicon
(royalty, assignment, patent prosecution, infringement, licensing, trade secret,
patentability, prior art, intellectual property, copyright, trademark, know-how,
commercialization, technology transfer, disclosure, conception, reduction to practice,
and joint invention).

For each sentence $i$ in policy $j$, we compute counts of supportive, restrictive,
sanction, legalese, and \IP-technical terms; the presence of second-person constructions
(\textit{you, your, yours}) and inclusive first-person constructions (\textit{we, our,
us}); the obligation-to-permission modal ratio; sentence length in words; and binary
indicators for very long sentences (exceeding 40 words) and for presence of procedural
cues.  Sentence-level features are aggregated to the policy level and normalized per
1,000 words to control for document length, then standardized ($z$-scored) using the
cross-policy mean and standard deviation.  The sub-index formulas are:
\begin{align*}
\text{ToneIndex}_j &= \tfrac{1}{6}\bigl[z(\text{supportive})_j + z(\text{2nd person})_j
  + z(\text{incl.\ we})_j - z(\text{restrictive})_j
  - z(\text{sanction})_j - z(\text{modal ratio})_j\bigr], \\[4pt]
\text{ClarityIndex}_j &= \tfrac{1}{3}\bigl[-z(\text{mean sent. length})_j
  - z(\text{long sent. share})_j + z(\text{procedural share})_j\bigr], \\[4pt]
\text{LegalLoadIndex}_j &= \tfrac{1}{2}\bigl[z(\text{legalese})_j
  + z(\text{IP-technical})_j\bigr].
\end{align*}

%% ── Appendix E ───────────────────────────────────────────────────────────────
\section{Robustness Checks}
\label{app:robust}

Cross-sectional results in Table~\ref{tab:crosssec} are robust to using the median
or trimmed mean rather than the mean for institution-level \PCI{} construction.  All
Kruskal-Wallis and Wilcoxon rank-sum tests yield the same direction and similar
$p$-values under all three aggregations.  The private R1 versus public R1 gap and the
land-grant gap are the most stable results; the urbanicity difference remains small
and statistically marginal under all aggregations.  Approximately 18 percent of
documents in the corpus required OCR processing.  Mean \PCI{} for OCR-processed
documents is not statistically distinguishable from the mean for clean documents, and
all cross-sectional results are unchanged when OCR-processed documents are excluded.
Restricting the sample to primary \IP{} policies only also yields results that are
qualitatively unchanged: the private R1 versus public R1 gap and the land-grant gap
remain directionally consistent and of comparable magnitude.

As a falsification exercise, we randomly permute institution labels 1,000 times and
compare observed cross-sectional gaps to the resulting null distributions.  The
observed private R1 versus public R1 \PCI{} gap is exceeded by permuted gaps in fewer
than 5 percent of cases.  This confirms that the documented cross-sectional patterns
reflect genuine signal in policy language rather than artifacts of the measurement
pipeline.  Table~\ref{tab:perm} reports the full permutation test results.

\begin{table}[H]
\centering
\caption{Placebo Permutation Tests}
\label{tab:perm}
\begin{threeparttable}
\small
\begin{tabular}{lccc}
\toprule
\textbf{Comparison} & \textbf{Observed Gap} & \textbf{Permutation $p$-value} & \textbf{$N$ Permutations} \\
\midrule
Private R1 vs.\ Public R1     & 0.030 & $<$0.05 & 1,000 \\
Med School vs.\ No Med School & 0.017 & $<$0.30 & 1,000 \\
\bottomrule
\end{tabular}
\begin{tablenotes}
\small
\item \textit{Notes:} Two-sided permutation $p$-value.  Exact values reported in
  the replication package.
\end{tablenotes}
\end{threeparttable}
\end{table}

%% ── Appendix F ───────────────────────────────────────────────────────────────
\section{Institution-Level PCI Reference Table}
\label{app:ref}

Full institution-level \PCI{} scores, comprising the mean across all observed policy
versions, the most recently observed \PCI{} score, the most recent year in the corpus,
the number of policies, and the standardized sub-indices, are provided in the file
\texttt{appendix\_F\_institution\_pci.csv} in the replication package.
Figure~\ref{fig:ranked} provides a visual summary for all 141 institutions.

\begin{figure}[H]
\centering
\includegraphics[width=0.78\textwidth]{fig8_ranked_institutions.pdf}
\caption{Institution-Level PCI Scores, All Research Universities}
\label{fig:ranked}
\floatfoot{\textit{Notes:} All 141 institutions sorted by mean \PCI.
  Color indicates institutional type.  Dashed line: neutral midpoint (0.50).
  Dash-dot line: overall mean.}
\end{figure}


\input{ranking}

\end{appendices}

\end{document}